% -------------------------------------------------------
% ASSIGNED TO: \textcolor{green}{Person 5}
% -------------------------------------------------------
% INSTRUCTION: Paste existing Section 4.6 Experiments here.
% ADD the CCLE robustness validation paragraph (marked in red below).
% Convert first-person to third person throughout.
% -------------------------------------------------------

% [Paste Section 4.6 Experiments here]
The research follows a two-stage experimental design to isolate the impact of model sharpness on explanation quality.
In Stage 1, the Overfitting Baseline is established as the RED XAI diagnostic starting point. DeepCCDS is trained using the Adam 
optimiser until a stable generalisation gap has emerged using the early stopping criterion defined in Section 4.1.4, confirming the 
Clever Hans baseline. Rules are then extracted from this overfitted model by both the GA and NSGA-II pipelines. The resulting rule sets 
are audited using PROMETHEE II ranking and GSEA pathway enrichment analysis to document the Fragility Baseline scorecard.
In Stage 2, the Regularisation Intervention is applied. The same DeepCCDS architecture is retrained using the SAM optimiser to smooth the 
loss landscape and reduce the generalisation gap, as described in Section 4.4. Rules are re-extracted using identical GA and NSGA-II 
parameters to ensure comparability. The same PROMETHEE II and GSEA audit is performed to generate a second evaluation scorecard. 
Wilcoxon signed-rank tests are then applied to the distributions of PROMETHEE outranking flows from both stages to determine whether 
the improvement in rule quality is statistically significant.


% -------------------------------------------------------
% NEW CONTENT — SECOND DATASET FOR ROBUSTNESS (supervisor suggestion)
% -------------------------------------------------------

\textcolor{red}{
% ADD: CCLE as a planned robustness validation dataset.
% The supervisor suggested this if time permits. Frame it as a stretch objective.
% Suggested text:
\paragraph{Robustness Validation: Cancer Cell Line Encyclopedia (CCLE)}
As a stretch objective, subject to timeline constraints, the complete RED XAI pipeline will be
applied to a second pharmacogenomic dataset: the Cancer Cell Line Encyclopedia (CCLE)
\cite{Barretina2012}. The CCLE provides gene expression profiles and drug response data for
cancer cell lines using a structurally compatible format to GDSC, and has been used alongside
GDSC in several of the studies cited in this proposal \cite{Partin2023, Baptista2021}. Applying
the pipeline to CCLE data, with LN IC\textsubscript{50} as the response target, would enable
assessment of whether the observed advantage of SAM regularisation over the Adam baseline
in terms of rule fidelity, comprehensibility, and biological validity generalises beyond the GDSC
dataset. Consistent results across both datasets would substantially strengthen the robustness
of the study's conclusions. No new statistical machinery is required; the same Wilcoxon
signed-rank test and PROMETHEE II framework are applied, and results are reported
comparatively.
% NOTE: Barretina et al. (2012) "The Cancer Cell Line Encyclopedia enables predictive modelling
% of anticancer drug sensitivity." Nature, 483, 603-607. Add to references.bib.
}